\begin{ExerciseBlock}
\ExerciseHeading{03}{Komet}{angular_momentum_comet_orbit}
Der Halleysche Komet, der erst im Jahr 2061 wieder am Himmel sichtbar sein wird, bewegt sich auf einer stark elliptischen Bahn um die Sonne (siehe Abbildung 9). Der Ort des kleinsten Abstands zur Sonne heißt Perihel (Punkt P) und der Ort des größten Abstandes Aphel (Punkt A). Im Perihel hat der Halleysche Komet einen Abstand von $r_P = 8,77 \cdot 10^7 \,$km zur Sonne und erreicht eine Geschwindigkeit von $v_P = 54,6 \,$km/s. Im Aphel hat der Komet einen Abstand von $r_A = 5,25 \cdot 10^9\,$km. Berechnen Sie die Geschwindigkeit $v_A$ des Kometen im Aphel. [Hinweis: Im Aphel und Perihel sind Geschwindigkeits- und Ortsvektor (bezüglich der Sonne) senkrecht zueinander.]

\begin{center}
	\CachedTikzFigure{content/resources/chapters/angular_momentum/tikz/comet_orbit.tex}
\end{center}
\end{ExerciseBlock}
\UseExerciseSolution{angular_momentum_comet_orbit}
